% Draft 1 (2026-11-16). Built from notes/2026-11-15_synthesis_steps23-28.md and the per-step notes.
% Every number below is taken from a step note; the step is given in a comment next to it.
\documentclass[aps,pra,twocolumn,superscriptaddress,nofootinbib]{revtex4-2}
\usepackage{amsmath,amssymb,bm}
\usepackage{graphicx}

\newcommand{\Gm}{G}
\newcommand{\ME}{\mathrm{MaxEnt}(3)}

\begin{document}

% Title: option A, chosen by the author on 2026-11-16.
\title{Completing high-order correlations of noisy IQP circuits with a connectivity prior}

\author{Marcus Tseng}
\altaffiliation{Undergraduate in the CYCU--UWM Dual Degree Program in Electrical Engineering and Computer Engineering.}
\affiliation{College of Electrical Engineering and Computer Science, Chung Yuan Christian University, Taoyuan 320314, Taiwan}

\date{\today}

\begin{abstract}
A quantum device's samples determine its low-order correlations well, but its high-order correlations poorly.
We ask whether a quantum circuit model can complete the high-order parity moments better than the classical model
that sees the same information, the maximum-entropy model on order-three moments.
A generic IQP circuit cannot, because it is not identifiable from these moments.
We show that the device's coupling graph fixes this.
On sixteen-qubit grid circuits, an IQP restricted to the graph completes orders four to six better than the
maximum-entropy model on all fourteen fresh instances, with a median error ratio of 0.35.
A recipe of learned per-qubit damping, a residual alarm and graph enlargement keeps this advantage under readout noise,
missing couplers and independent or correlated gate noise up to two percent.
Coherent over-rotations of about one percent per gate are the limit, and a coherent-aware model did not remove it.
\end{abstract}

\maketitle

% ---------------------------------------------------------------------------
\section{Introduction}

Quantum generative models are often proposed as a route to a practical quantum advantage
\cite{huang2025generative,hibatallah2024framework}, but the evidence so far is mixed.
IQP circuit Born machines can be trained classically at large scale, yet their reported wins are against classical
models that were not given the same information \cite{recio2025train}.
The authors of that study note that their classical baselines could likely be improved.
A related study reports a large win over a maximum-entropy model fitted to parity moments \cite{baumann2026parity}.
That classical model had no order cutoff and was fitted to two hundred samples, so it may have overfitted.

We use a fairer test, in which both models receive the same order-$\le 3$ parity moments estimated from $N$ shots.
Each model must then predict the order-4 to order-6 moments, which the shots determine poorly.
The classical opponent is $\ME$, the maximum-entropy distribution with the given moments \cite{jaynes1957information}.
It is the least-committal completion of that information, so a model beats it only if its inductive bias matches
the structure of the data.

Our first finding is negative: a generic IQP circuit with all 696 order-$\le 3$ generators is not identifiable from
these moments, and it loses to $\ME$ on most instances.
Our second finding is positive.
Restricting the IQP to the device's coupling graph makes it identifiable, and the restricted model beats $\ME$
as well as a maximum-entropy model given the same graph.
Our third finding is practical: a three-step recipe keeps the advantage under realistic departures from the ideal
circuit, and we locate the noise class where it fails.

Every result in this paper comes from simulation of sixteen qubits on a $4\times 4$ grid.
Each robustness claim was tested on fresh random instances under a protocol fixed before the run.
The claim is statistical, not computational.
Sampling from IQP circuits is believed to be classically hard \cite{bremner2011classical}, but their parity moments
can be estimated classically at any size \cite{recio2025train}.

% ---------------------------------------------------------------------------
\section{Setup}

\subsection{Task and score}

Let $p$ be a distribution on $n=16$ bits.
The parity moment of a mask $S\subseteq\{1,\dots,n\}$ is $f_S = \mathbb{E}_{x\sim p}[(-1)^{\sum_{j\in S} x_j}]$.
A model receives the 696 estimated moments $\hat f_S$ with $1\le |S|\le 3$, computed from $N=10^4$ shots, and
outputs a distribution $q$.
We score its completion of orders four to six against the exact target:
\begin{equation}
  \Gm(q) = \frac{1}{3}\sum_{l=4}^{6} \frac{1}{\binom{n}{l}} \sum_{|S|=l} \bigl(f_S(q) - f_S(p)\bigr)^2 .
\end{equation}
The main statistic is the ratio $R = \Gm(\text{model})/\Gm(\ME)$, and $R<1$ means that the model completes better
than $\ME$.

\subsection{Models}

An IQP circuit with generator set $\mathcal{G}$ prepares
$|\psi\rangle = H^{\otimes n}\exp(i\sum_{g\in\mathcal{G}}\theta_g Z_g)H^{\otimes n}|0\rangle$,
where $Z_g$ is the product of Pauli $Z$ operators on the qubits of mask $g$.
Its parity moments have a closed form, and we fit $\theta$ by least squares on the order-$\le 3$ moments.
The fit uses L-BFGS-B with twenty random restarts and keeps the solution with the lowest residual.

The \emph{local IQP} uses only the generators allowed by the coupling graph.
On the $4\times4$ grid these are the 16 single-qubit terms, the 24 edges, the 36 corner triples of the 9 plaquettes, and
the 9 plaquettes, which gives 85 candidates in all.

\emph{FlipIQP} adds one damping factor per qubit to model noise.
It multiplies each output moment by $\prod_{j\in S}\lambda_j$, with $\lambda_j\in[0.5,1]$.
These 16 factors absorb readout and gate noise without any calibration data.

The classical arms are $\ME$, with 696 parameters, and a maximum-entropy model restricted to the 76 local masks of
weight at most three.
We call the second model MaxEnt-on-graph, since it receives the same prior as the local IQP.

\subsection{Data}

The target is an IQP circuit on the grid with 49 generators: the singles, the edges and the plaquettes.
The model is not told which of its candidates are active.
Angles are drawn as $\theta_g\sim\mathcal{N}(0,\sigma^2)$ with $\sigma=0.251$.
Later sections add readout bit flips, Pauli errors after each gate, or coherent over-rotations.

\subsection{Protocol}

Each test was specified in writing before it was run, including its hypotheses, thresholds and fresh random seeds.
A hypothesis passed if the model beat its opponent on at least 11 of 14 instances, a one-sided sign test at $p=0.029$.
Missing or failed runs counted as losses.
Development used separate seeds, which were never reused for testing.

% ---------------------------------------------------------------------------
\section{Why a prior is needed}

The order-$\le3$ moments do not determine the higher moments.
For any target $p$, every mixture of $p$ with its own $\ME$ distribution has the same order-$\le3$ moments.
The high-order moments of these mixtures, however, interpolate between those of $p$ and those of $\ME$.
No method that sees only low-order moments can be right on every such mixture, so an advantage must come from a
prior that matches the data.

The generic IQP has no such prior.
With all 696 order-$\le3$ generators its fit is badly conditioned, with a condition number near $10^9$.
At $N=10^4$ it beats $\ME$ on only 2 of 14 instances, with median $R=1.17$.

The coupling graph supplies the missing prior.
Table~\ref{tab:step23} shows that the local IQP beats $\ME$ on all 14 fresh targets, at both $N=10^4$ and $N=10^5$.
The advantage grows with $N$, as expected when the model family is right and $\ME$ is not.

\begin{table}[t]
\caption{Local IQP on fresh grid targets. $R$ is the median ratio of completion errors; wins are out of 14.}
\label{tab:step23}
\begin{ruledtabular}
\begin{tabular}{lccc}
Opponent & $N$ & Wins & Median $R$ \\
\hline
$\ME$ & $10^4$ & 14 & 0.35 \\
$\ME$ & $10^5$ & 14 & 0.19 \\
MaxEnt-on-graph & $10^4$ & 14 & 0.068 \\
Generic IQP & $10^4$ & 14 & --- \\
\end{tabular}
\end{ruledtabular}
\end{table}
% step 23: H1-H4, seeds 920-933.

The classical model cannot use the same prior.
MaxEnt-on-graph has about as many parameters as the local IQP, 76 against 85, yet it completes far worse.
Interference makes $\log p$ non-local in the graph, so a local log-linear model is the wrong family for these data.

The prior is not a general improvement.
On a classical grid-Ising target the local IQP loses to $\ME$ on all 14 instances.
An IQP model also gave no advantage on real surface-code detection events \cite{google2023suppressing}.
The advantage appears only when the data come from a circuit like the model.

% ---------------------------------------------------------------------------
\section{Robustness}

A real device departs from the ideal circuit in two ways: its stated coupling graph may be wrong, and its gates and
readout are noisy.
Table~\ref{tab:robust} summarizes the tests of both departures.
Each row is a separate test on fresh instances.

\begin{table*}[t]
\caption{Robustness tests on the $4\times4$ grid. Each row uses 14 fresh instances, $N=10^4$, and compares with
$\ME$. ``Wrong graph'' means two active couplers are missing from the model. Rows marked $^\dagger$ were run in the same
protocol but were reported without a pre-specified hypothesis.}
\label{tab:robust}
\begin{ruledtabular}
\begin{tabular}{llcc}
Departure & Method & Wins & Median $R$ \\
\hline
2 missing couplers & local IQP & 3 & 2.86 \\
18 extra inactive candidates & local IQP & 14 & 0.46 \\
2\% readout noise & local IQP & 11 & 0.78 \\
2\% readout noise & local IQP, mitigated$^\dagger$ & 14 & 0.42 \\
Gate noise 0.5\% & FlipIQP & 14 & 0.51 \\
Gate noise 1\% & FlipIQP & 13 & 0.63 \\
Gate noise 2\% & FlipIQP$^\dagger$ & 13 & 0.60 \\
Wrong graph, gate noise 0.5\% & recipe & 13 & 0.56 \\
Wrong graph, gate noise 1\% & recipe & 12 & 0.66 \\
Correct graph, gate noise 1\% & recipe & 13 & 0.57 \\
Wrong graph, correlated noise 1\% & recipe & 14 & 0.54 \\
Correct graph, correlated noise 1\% & recipe & 14 & 0.53 \\
Wrong graph, correlated noise 2\% & recipe & 12 & 0.80 \\
\end{tabular}
\end{ruledtabular}
\end{table*}
% rows: step 24 (M1, E1, D1, mitigated arm reported), step 26 (G2, G1, stress), step 27 (K2, K1, K4), step 28 (C1, C2, C4).
% Gate-noise rows include 2% readout noise.

\subsection{Graph errors}

Extra candidates are cheap: adding 18 inactive diagonal candidates moves the median $R$ only from 0.35 to 0.46.
Missing candidates are costly, and dropping two of the 24 couplers turns the win into a clear loss.
The prior must therefore contain every active coupler.

A missing coupler leaves a mark in the fit, because the training residual can no longer reach the sampling-noise
floor $F=\sum_S(1-\hat f_S^2)/N$.
We raise an alarm when the residual exceeds $0.75F$.
When the alarm fires, the model is refitted on an enlarged graph that adds the 18 plaquette diagonals, for 103
candidates in all.

\subsection{Noise}

Readout noise alone can be mitigated if its rate is calibrated, but gate noise cannot be removed this way.
With readout mitigation only, the advantage fades at 1\% gate error and reverses at 2\%.
FlipIQP's learned damping restores it, keeping the median $R$ near 0.6 from 0.5\% to 2\% gate error.
The learned factors track the noise, and their median falls from 0.94 to 0.86 as the gate error grows.

\subsection{The recipe}

The final recipe has three steps.
(i) Fit FlipIQP on the stated coupling graph, using the raw order-$\le3$ moments.
(ii) If the residual exceeds $0.75F$, refit on the enlarged graph.
(iii) Use the fitted model to predict the higher moments.
The recipe needs no calibration of readout or gate errors.

The recipe passed every one of its tests in Table~\ref{tab:robust}.
Without the alarm, a wrong graph loses to $\ME$ on most instances, winning only 6 of 14 at 1\% gate noise.
On the correct graph the alarm fired once in 14 instances at 1\% independent noise, and never at 1\% correlated noise.
FlipIQP assumes noise that is independent across qubits.
Correlated multi-qubit Pauli errors break that assumption, yet the recipe still won all 14 instances at 1\%.

The alarm weakens as the noise grows.
At 2\% correlated noise, four of the 14 wrong-graph instances stayed below the threshold, and two of these lost.
Always enlarging the graph won all 14 instances in that setting, with median $R=0.72$.
On a correct graph enlarging costs some accuracy, raising the median $R$ from 0.57 to 0.71 at 1\% noise.
We therefore recommend always enlarging when gate errors reach about 2\%, a recommendation we observed but did not
test.

% ---------------------------------------------------------------------------
\section{The limit: coherent over-rotations}

Pauli noise shrinks correlations, which FlipIQP's damping can model, but a coherent error does not shrink them.
We tested over-rotations $R_X(\epsilon_j)$ applied to qubit $j$ after every gate that touches it.
Each $\epsilon_j$ is a fixed miscalibration drawn from $\mathcal{N}(0,\epsilon^2)$.
The infidelity per gate is about $\epsilon^2/4$, so $\epsilon=0.2$ corresponds to about 1\%.

On a wrong graph the recipe still passed, winning 12 of 14 instances at $\epsilon=0.2$ and 13 of 14 at
$\epsilon=0.1$.
On the correct graph it failed at $\epsilon=0.2$, winning only 10 of 14 with median $R=0.74$.
The alarm also fired falsely on 3 of those 14 instances.
The learned damping stayed near $\lambda=1$ and therefore did nothing.
A unitary error moves weight between moments rather than shrinking them, so damping is the wrong correction.

The failure is not visible in the fit.
On one losing instance the residual was only $0.35F$, yet $R=1.21$.
The residual checks orders up to three, while the damage appears at higher orders.

We then explored a coherent-aware model that places unknown $R_X$ layers between the gate layers of the IQP circuit.
This exploration reused instances from the failed test, so it has no formal verdict.
The new model fitted the order-$\le3$ moments better than FlipIQP but did not complete the higher orders better.
It beat FlipIQP on only 7 of 14 coherent instances, and on only 2 of 14 instances with Pauli noise, where it overfitted.
A version with one angle per qubit, which matches how the data were generated, did no better.
At $N=10^4$ the low-order moments do not pin down the coherent angles.

% ---------------------------------------------------------------------------
\section{Discussion}

A prior from the device's coupling graph turns an unidentifiable IQP model into a useful one.
The resulting model completes high-order moments better than the information-matched classical model.
Without noise, it also beats the classical model given the same graph by more than an order of magnitude.
The advantage survives the departures a real device is likely to show, with the exception of coherent over-rotations.

The scope of the claim is narrow.
It covers data from IQP-like devices with an approximately known coupling graph.
It does not cover classical data, IQP models gave no advantage on surface-code detection events, and the advantage
is statistical rather than computational.

Three questions remain open.
The first is coherent miscalibration, which more shots or calibrated rotation angles might fix, although neither was tested.
The second is size, because the exact score needs $2^n$ probabilities and larger tests need an estimated score or a
cheaper fitter.
The third is real hardware data, which is the test that matters most.

% ---------------------------------------------------------------------------
\begin{acknowledgments}
% TODO: acknowledgments; ask whether an advisor should be an author or acknowledged.
\end{acknowledgments}

\bibliography{refs}

\end{document}
